%% ============================================================
%\part{گاف انرژی: برآورد، مهندسی، پژوهش}
%% ============================================================
%
%\chapter{گاف انرژی در شبیه‌سازی کوانتومی}
%\label{chap:gap-energy}
%
%\section{جایگاه در نقشه}
%
%این فصل یک موضوع پژوهشی است، نه فقط یک کمیت کمکی برای فصل~\ref{chap:evo-adi}.
%در نقشه شکل~\ref{fig:research-map} گاف در چند جا ظاهر می‌شود: به‌عنوان پرسش طیفی، به‌عنوان گلوگاه مسیر آدیاباتیک، و به‌عنوان ویژگی فیزیکی خودِ مدل.
%بخش‌های~\ref{sec:gap-research}--\ref{sec:gap-projects} گاف را در چارچوب یک مسیر $H(s)$ و مثال دوکیوبیتی باز می‌کنند.
%اینجا همان کمیت را از آن مسیر جدا می‌کنیم و می‌پرسیم: در شبیه‌سازی کوانتومی چرا گاف مهم است، چگونه آن را می‌یابند یا بزرگ می‌کنند، و کدام مسئله‌های باز می‌توانند به مقاله منجر شوند.
%
%هدف عملی فصل این است که پژوهش جاری گروه روی گاف متمرکز شود: مسئله، روش، معیار موفقیت، و رقیب کلاسیک از ابتدا روشن باشند.
%
%\section{ایده شهودی}
%
%گاف انرژی فاصله میان دو بخش طیف است؛ معمولاً فاصله حالت پایه از نخستین حالت برانگیخته مرتبط.
%این عدد کوچک به نظر می‌رسد، ولی نقش‌های کاملاً متفاوتی بازی می‌کند:
%\begin{itemize}
%	\item در محاسبه آدیاباتیک، گاف کمینه مسیر زمان اجرا را کنترل می‌کند (بخش~\ref{sec:gap-runtime}).
%	\item در فیزیک بس‌ذره‌ای، گاف می‌گوید سامانه عایق است یا فلز، ابررسانا است یا بهنجار، توپولوژیک است یا بدیهی.
%	\item در برآورد فاز، تفکیک دو ویژه‌مقدار به گاف وابسته است.
%	\item در سامانه باز و حافظه کوانتومی، گاف مقیاس حفاظت در برابر تحریک گرمایی و برخی خطاهاست.
%\end{itemize}
%بنابراین «یافتن گاف» و «بزرگ کردن گاف» دو مسئله جدا هستند که هر دو زیرمجموعه شبیه‌سازی‌اند.
%اولی یک \emph{پرسش طیفی} است؛ دومی یک \emph{مهندسی مسیر یا رمزگذاری} است.
%شکل~\ref{fig:gap-methods-map} این تفکیک را نشان می‌دهد.
%
%\begin{figure}[htbp]
%	\centering
%	\begin{latin}
%		\setstretch{1}
%		\resizebox{\ifdim\width>\textwidth\textwidth\else\width\fi}{!}{%
%			% دسته‌بندی روش‌های برآورد و مهندسی گاف
% سبک مشترک نمودارهای سند پایه‌ها (فقط داخل محیط latin فراخوانی شود)
\tikzset{
	qafbox/.style={
		draw=black!55,
		rounded corners=2pt,
		align=center,
		inner sep=3pt,
		thick,
		fill=white,
		font=\footnotesize\linespread{1}\selectfont,
		minimum height=0.95cm
	},
	qafwide/.style={
		qafbox,
		text width=2.55cm
	},
	qafnarr/.style={
		qafbox,
		text width=2.05cm
	},
	qaftitle/.style={
		font=\small\bfseries\linespread{1}\selectfont,
		align=center
	},
	qafarr/.style={
		-{Latex[length=2.2mm]},
		thick,
		draw=black!65
	},
	qafpanel/.style={
		draw,
		thick,
		rounded corners=4pt,
		inner xsep=8pt,
		inner ysep=8pt
	}
}

\begin{tikzpicture}[
	x=1cm,
	y=1cm,
	font=\footnotesize,
	gbox/.style={qafbox, text width=2.35cm, minimum height=1.05cm}
	]
	\node[gbox, fill=blue!12, text width=3.4cm] (why) at (0,3.4) {Why the gap?};
	\node[gbox, fill=green!12] (cl) at (-5.2,1.6) {Classical spectrum};
	\node[gbox, fill=orange!14] (dyn) at (-1.75,1.6) {Dynamics / spectroscopy};
	\node[gbox, fill=cyan!14] (var) at (1.75,1.6) {Variational / subspace};
	\node[gbox, fill=red!12] (ft) at (5.2,1.6) {FTQC / QPE};
	\node[gbox, fill=green!10, text width=2.55cm] (cl2) at (-5.2,0) {Lanczos, DMRG, QMC};
	\node[gbox, fill=orange!10, text width=2.55cm] (dyn2) at (-1.75,0) {Oscillation, quench, $C(\tau)$};
	\node[gbox, fill=cyan!10, text width=2.55cm] (var2) at (1.75,0) {VQD, QSE, Krylov};
	\node[gbox, fill=red!10, text width=2.55cm] (ft2) at (5.2,0) {Phase est., counting};
	\node[gbox, fill=violet!12, text width=4.6cm] (eng) at (0,-1.7) {Gap engineering / shortcuts};
	\node[font=\scriptsize, text=black!60] at (0,-2.45) {schedule, catalyst, CD, encoding};
	\draw[qafarr] (why) -- (cl);
	\draw[qafarr] (why) -- (dyn);
	\draw[qafarr] (why) -- (var);
	\draw[qafarr] (why) -- (ft);
	\draw[qafarr] (cl) -- (cl2);
	\draw[qafarr] (dyn) -- (dyn2);
	\draw[qafarr] (var) -- (var2);
	\draw[qafarr] (ft) -- (ft2);
	\draw[qafarr] (cl2.south) -- ++(0,-0.35) -| (eng.north);
	\draw[qafarr] (dyn2.south) -- ++(0,-0.35) -| (eng.north);
	\draw[qafarr] (var2.south) -- ++(0,-0.35) -| (eng.north);
	\draw[qafarr] (ft2.south) -- ++(0,-0.35) -| (eng.north);
\end{tikzpicture}
%
%		}
%	\end{latin}
%	\caption{چهار خانواده برآورد گاف و یک خانواده مهندسی. برآورد می‌گوید گاف چیست؛ مهندسی می‌گوید مسیر یا کنترل را چگونه عوض کنیم.}
%	\label{fig:gap-methods-map}
%\end{figure}
%
%\section{چرا گاف در این حوزه مهم است؟}
%
%\subsection{چند معنای فیزیکی که نباید قاطی شوند}
%
%کلمه «گاف» در مقالات مختلف یک کمیت واحد نیست.
%برای پژوهش قابل‌انتشار باید از ابتدا مشخص شود کدام گاف مد نظر است~\cite{Georgescu2014,AlbashLidar2018}.
%\begin{enumerate}
%	\item \textbf{گاف بس‌ذره‌ای کم‌انرژی.}
%	$\Delta=E_1-E_0$ برای هامیلتونی بسته با تعداد ذره یا تقارن ثابت.
%	همین تعریف در فصل~\ref{chap:evo-adi} برای مسیر آدیاباتیک به کار رفت.
%	\item \textbf{گاف بخش تقارنی.}
%	اگر $[H,S]=0$، گذار میان دو بخش ممکن است ممنوع باشد؛ گاف کل می‌تواند صفر شود در حالی که گاف بخش قابل‌دسترس مثبت بماند (بخش~\ref{sec:gap-research}).
%	\item \textbf{گاف بار / گاف اسپینی / گاف اپتیکی.}
%	در مدل هابارد یا مولکول، فاصله دو ویژه‌مقدار کم‌انرژی الزاماً همان انرژی لازم برای افزودن یک الکترون یا جذب یک فوتون نیست.
%	قواعد انتخاب مشاهده‌پذیر تعیین می‌کند کدام گاف در آزمایش دیده می‌شود.
%	\item \textbf{گاف توپولوژیک.}
%	در فازهای با ناورداهای توپولوژیک، گاف حجمی سامانه را از اختلال موضعی محافظت می‌کند؛ حالت‌های لبه ممکن است گاف کوچک یا صفر داشته باشند.
%	\item \textbf{گاف لیوویلی.}
%	در سامانه باز، نرخ واهلش را گاف مولد دینامیک باز می‌دهد، نه لزوماً $E_1-E_0$ی هامیلتونی بسته~\cite{GapSels2017}.
%\end{enumerate}
%اگر مقاله ما این‌ها را یکی بگیرد، نتیجه برای داور قابل تفسیر نیست.
%
%\subsection{نقش گاف در محاسبه آدیاباتیک}
%
%در رهیافت آدیاباتیک، تخمین مقیاسی $T\sim\hbar/\Delta_{\min}^{2}$ گاف را به منبع تبدیل می‌کند~\cite{Farhi2000,GapJansen2007,AlbashLidar2018}.
%بسته‌شدن چندجمله‌ای در برابر نمایی همان تمایزی است که بخش~\ref{sec:complexity-scaling} باز کرد.
%پژوهش گاف در این شاخه سه سؤال جدا دارد:
%\begin{itemize}
%	\item گاف لحظه‌ای $\Delta(s)$ در یک مدل مشخص چیست؟
%	\item کمینه آن روی مسیر کجاست و با $n$ چگونه بسته می‌شود؟
%	\item آیا با عوض کردن مسیر، زمان‌بندی، یا جمله کمکی می‌توان $\Delta_{\min}$ را بزرگ‌تر کرد، یا گذار را بدون بزرگ کردن گاف سرکوب کرد؟
%\end{itemize}
%سؤال سوم دیگر «محاسبه گاف» نیست؛ مهندسی مسیر است و در بخش~\ref{sec:gap-eng} می‌آید.
%
%\subsection{نقش گاف در شبیه‌سازی دینامیکی و مواد}
%
%در شبیه‌سازی تحول زمانی، گاف مقیاس نوسان مشاهده‌پذیرها و آستانه پاسخ خطی را می‌دهد~\cite{Georgescu2014,Daley2022}.
%اگر گاف کوچک باشد، تفکیک طیفی نیاز به زمان همدوس بلند دارد: $\delta E\sim 2\pi\hbar/t_{\max}$.
%در شیمی کوانتومی، گاف هومو--لومو یا شکافت سینگلت--تریپلت کمیت فیزیکی مستقلی است، نه فقط گلوگاه الگوریتم~\cite{McArdle2020,Bauer2020}.
%در شبیه‌ساز آنالوگِ شبکه نوری یا اتم خنثی، طیف‌سنجی پس از quench یکی از خروجی‌های اصلی آزمایش است؛ گاف از مکان قله ساختار دینامیکی خوانده می‌شود، نه از قطری‌سازی کلاسیکِ ماتریس کامل.
%
%\subsection{نقش گاف در کیوبیت و سخت‌افزار}
%
%برای کیوبیت ابررسانا، نایکسانی ترازها همان چیزی است که دو تراز پایین را از بقیه جدا می‌کند (فصل~\ref{chap:models}).
%گاف بزرگ‌تر به معنی حفاظت بهتر در برابر تحریک گرمایی است، به شرط آن‌که عملگر جفت‌شدگی محیط با آن گذار هم‌پوشانی داشته باشد.
%در کدهای پایدارکننده، هامیلتونی جریمه خطا گاف ایجاد می‌کند؛ آستانه تصحیح خطا به این گاف و به وزن خطا وابسته است، نه فقط به تعداد کیوبیت.
%پس پژوهش گاف فقط «الگوریتم آدیاباتیک» نیست: طراحی کیوبیت، حافظه، و شبیه‌سازی آنالوگ هم مصرف‌کننده همان مفهوم‌اند.
%
%\section{رهیافت‌های اساسی برای یافتن گاف}
%
%جدول~\ref{tab:gap-method-families} چهار خانواده را با هزینه، خروجی، و رژیم نوعی جمع می‌کند.
%جزئیات عددی مسیر دوکیوبیتی در بخش~\ref{sec:gap-methods} آمده است؛ اینجا هر خانواده را به‌عنوان خط پژوهشی باز می‌کنیم.
%
%\begin{table}[htbp]
%	\centering
%	\small
%	\setstretch{1.12}
%	\caption{چهار خانواده برآورد گاف در شبیه‌سازی کوانتومی. هزینه نوعی است، نه کران قضیه.}
%	\label{tab:gap-method-families}
%	\begin{tabularx}{\textwidth}{@{}
%		>{\raggedright\arraybackslash}p{3.3cm}
%		>{\raggedright\arraybackslash}X
%		>{\raggedright\arraybackslash}p{3.4cm}
%		@{}}
%		\toprule
%		\rowcolor{gray!14}
%		\textbf{خانواده} & \textbf{ایده} & \textbf{رژیم نوعی} \\
%		\midrule
%		طیف کلاسیک
%		& قطری‌سازی، لانچوس، \lr{DMRG}، \lr{QMC}
%		& $n$ کوچک یا درهم‌تنیدگی محدود
%		\\
%		\rowcolor{gray!8}
%		طیف‌سنجی دینامیکی
%		& نوسان $\langle O(t)\rangle$، quench، همبستگی زمان موهومی
%		& آنالوگ و دیجیتال؛ وابسته به مشاهده‌پذیر
%		\\
%		وردشی / زیرفضا
%		& \lr{VQD}، \lr{SSVQE}، \lr{QSE}، کریلوف کوانتومی
%		& \lr{NISQ} و \lr{early-FTQC}
%		\\
%		\rowcolor{gray!8}
%		برآورد فاز / شمارش
%		& \lr{QPE}، \lr{QSVT}، پرسمان شمارش
%		& تحمل‌پذیر در برابر خطا
%		\\
%		\bottomrule
%	\end{tabularx}
%\end{table}
%
%\subsection{طیف کلاسیک: مرجع، نه رقیب همیشه مغلوب}
%
%برای هر مقاله روش کوانتومی، مرجع کلاسیک باید صریح باشد~\cite{GapSchollwock2011,Cirac2021}.
%قطری‌سازی کامل برای $D=2^n$ با حافظه $O(D^2)$ فقط تا حدود $n\simeq 12$--$16$ روی رایانه معمولی راحت است.
%لانچوس ضرب $Hv$ را بدون ماتریس متراکم انجام می‌دهد و چند تراز پایین را می‌دهد؛ همگرایی به گاف و به بردار شروع وابسته است.
%شبکه تانسوری (\lr{DMRG/MPS}) برای زنجیره کم‌درهم‌تنیده هنوز قوی‌ترین ابزار کلاسیک است.
%مونت‌کارلو کوانتومی در مدل بدون مشکل علامت می‌تواند از افت نمایی همبستگی زمان موهومی گاف را بخواند.
%
%\textbf{نتیجه برای پروژه.}
%اگر مدل ما زنجیره \lr{TFIM} یا مولکول کوچک است، مرجع کلاسیک غالباً دقیق‌تر و ارزان‌تر از مدار نوفه‌دار است.
%سهم کوانتومی وقتی معنا دارد که یا مدل در رژیم سخت کلاسیک باشد (دوبعدی، فرمیون با علامت، درهم‌تنیدگی حجمی)، یا خودِ \emph{پروتکل اندازه‌گیری گاف روی سخت‌افزار} موضوع مقاله باشد، نه فقط عدد گاف.
%
%\subsection{طیف‌سنجی دینامیکی: گاف به‌عنوان فرکانس}
%
%اگر حالت آزمایشی هم‌پوشانی با پایه و یک برانگیخته داشته باشد و مشاهده‌پذیر عنصر ماتریسی غیرصفر داشته باشد،
%\begin{equation}
%	\langle O(t)\rangle
%	=
%	A
%	+
%	2\,\mathrm{Re}\bigl[c_0^*c_1 O_{01}e^{-i\Delta t/\hbar}\bigr]
%	+\cdots
%	\label{eq:gap-osc-research}
%\end{equation}
%فرکانس نوسان همان گاف است (رابطه~\eqref{eq:gap-oscillation} در بخش~\ref{sec:gap-methods}).
%در شبیه‌ساز آنالوگ همین ایده به شکل quench و تبدیل فوریه تابع ساختار دینامیکی درمی‌آید~\cite{Georgescu2014,Daley2022}.
%در زمان موهومی، شیب $\log C_O(\tau)$ به کمترین گافِ دارای وزن میل می‌کند.
%
%محدودیت اصلی این خانواده انتخاب $O$ است: نبود قله دلیل نبود تراز نیست.
%تقارن می‌تواند گذار را خاموش کند؛ نویز خوانش می‌تواند قله را جابه‌جا یا پهن کند.
%مقالات \lr{Lee--Myers--Scarola} و \lr{Cugini et al.} همین خط را با پردازش سیگنال و آماده‌سازی آدیاباتیک دنبال می‌کنند~\cite{GapLee2025,GapCugini2025}.
%
%\subsection{وردشی و زیرفضا: دو انرژی، یک تفاضل خطرناک}
%
%\lr{VQE} یک کران بالا برای $E_0$ می‌دهد، نه گاف~\cite{Peruzzo2014,Tilly2022}.
%برای تراز برانگیخته باید قید اضافه‌ای گذاشت.
%در \lr{VQD}، جریمه همپوشانی با حالت‌های پایین‌تر تراز بعدی را هدف می‌گیرد~\cite{GapHiggott2019}:
%\begin{equation}
%	{\cal L}_k(\theta)
%	=
%	\langle\psi(\theta)|H|\psi(\theta)\rangle
%	+
%	\sum_{j<k}\mu_j\bigl|\langle\psi_j|\psi(\theta)\rangle\bigr|^2.
%\end{equation}
%\lr{SSVQE} چند حالت را هم‌زمان در یک زیرفضای پارامتری بهینه می‌کند.
%بسط زیرفضای کوانتومی (\lr{QSE}) از تحریک‌های موضعیِ یک مرجع، ماتریس‌های $A$ و $S$ می‌سازد و مسئله تعمیم‌یافته ویژه را کلاسیک حل می‌کند~\cite{McClean2017QSE}.
%کریلوف کوانتومی همان ایده را با توان‌ها یا تحول‌های زمانیِ یک حالت مرجع می‌سازد؛ نمایش سخت‌افزاری تا ده‌ها سایت گزارش شده است~\cite{KrylovNatComm2025,GapYu2025,Kirby2023Krylov}.
%
%\textbf{هشدار روش‌شناختی.}
%تفاضل دو انرژی وردشی کران تضمینی برای $\Delta$ نیست.
%خطای هر تراز می‌تواند هم‌علامت یا ناهم‌علامت باشد؛ گاف تخمینی ممکن است خوش‌بینانه شود.
%هر ادعای گاف در این خانواده باید باقیمانده $Hv-Ev$، همگرایی زیرفضا، و در صورت امکان گواهی فاصله تا طیف را گزارش کند.
%
%\subsection{برآورد فاز و الگوریتم‌های تحمل‌پذیر}
%
%اگر بتوان $U=e^{-iH\tau}$ را با دقت کنترل‌شده ساخت، \lr{QPE} ویژه‌مقدارها را به فاز ترجمه می‌کند~\cite{NielsenChuang2010,McArdle2020}.
%خطای مطلق هر انرژی $\epsilon_E$ به خطای گاف حداکثر $2\epsilon_E$ منجر می‌شود (بخش~\ref{sec:gap-methods}).
%زمان همدوس نوعی نسخه استاندارد مثل $O(\hbar/\epsilon_E)$ است؛ هزینه شبیه‌سازی $U$ جداست.
%الگوریتم‌های مبتنی بر بلوک‌کدگذاری و تبدیل تکین (\lr{QSVT}) می‌توانند گاف $k$ام و نقطه میانی آن را با تضمین پیچیدگی برآورد کنند؛ در مدل \lr{QRAM} هزینه هنوز می‌تواند نسبت به بعد ماتریس نمایی بماند~\cite{GapCarrera2026}.
%این خط برای افق \lr{FTQC} است، نه برای دستگاه نوفه‌دار ده‌کیوبیتی.
%
%\section{افزایش گاف و میان‌بر آدیاباتیک}
%\label{sec:gap-eng}
%
%«گاف را بزرگ کنیم» با «گذار را کم کنیم» یکی نیست.
%چهار سازوکار جدا در ادبیات وجود دارد~\cite{AlbashLidar2018,GapSels2017,Gueneau2025CD}.
%
%\subsection{زمان‌بندی: گاف را عوض نمی‌کند}
%
%زمان‌بندی غیرخطی $s(t)$ همان مسیر $H(s)$ را با سرعت متغیر طی می‌کند.
%گاف لحظه‌ای عوض نمی‌شود؛ فقط زمان بیشتری نزدیک کمینه خرج می‌شود~\cite{GapRoland2002}.
%این ارزان‌ترین مداخله است و باید خط پایه هر مقاله مهندسی باشد.
%
%\subsection{تغییر مسیر و هامیلتونی کاتالیزگر}
%
%جمله $f(s)H_{\mathrm{cat}}$ با $f(0)=f(1)=0$ مسیر را در فضای هامیلتونی‌ها خم می‌کند.
%هدف معمولاً اجتناب از گذار مرتبه اول یا افزایش $\Delta_{\min}$ است.
%بهای آن غیرمحلی‌شدن، افزایش نرم مشتق، و نیاز به کنترل اضافی است.
%بزرگ شدن گاف در یک $s$ کافی نیست؛ کمینه کل مسیر و عناصر گذار باید دوباره محاسبه شوند (بخش~\ref{sec:gap-control}).
%
%\subsection{رانش جبران‌کننده: میان‌بر، نه لزوماً گاف بزرگ‌تر}
%
%رانش جبران‌کننده\ft{counterdiabatic driving} گذار دیاباتیک را با یک جمله کمکی حذف می‌کند، حتی اگر گاف مرجع کوچک بماند~\cite{Berry2009STA,GapSels2017}.
%جمله دقیق معمولاً غیرمحلی است و خودِ طیف را می‌خواهد؛ بنابراین به‌صورت خام مسئله را دور می‌زند.
%تقریب‌های موضعی و بسط وردشیِ پتانسیل پیمانه‌ای آدیاباتیک عملی‌ترند، ولی نزدیک گاف نمایی کوچک اثرشان محدود می‌شود~\cite{Gueneau2025CD}.
%پیاده‌سازی دیجیتال، دیجیتال--آنالوگ روی یون، و آنالوگ روی اتم خنثی گزارش شده است~\cite{Kumar2024DACQO,ACQC2026}.
%برای گروه ما سؤال پژوهشی دقیق این است: در مدل اسپینی محلی با منابع کنترل محدود، تقریب موضعی \lr{CD} تا چه اندازه $\Delta_{\min}$ نمایی را جبران می‌کند، و هزینه گیت یا پالس در برابر جاروب آهسته‌تر چیست؟
%
%\subsection{رمزگذاری و تقارن}
%
%گاهی گاف کوچک مصنوع رمزگذاری است، نه فیزیک مسئله.
%شکستن یا حفظ تقارن، انتخاب هامیلتونی آغازین، و نگاشت \lr{QUBO} به سخت‌افزار همگی $\Delta_{\min}(n)$ را عوض می‌کنند.
%این شاخه به بهینه‌سازی ترکیبی نزدیک است، ولی اگر مدل فیزیکی اسپینی باشد همچنان شبیه‌سازی است: سؤال این است که کدام نمایش همان فیزیک را با گاف ملایم‌تر می‌دهد.
%
%\section{کاربرد در سامانه‌های کیوبیتی و شبیه‌سازی}
%
%چهار کاربرد مشخص با مسئله باز قابل‌تبدیل به مقاله:
%
%\paragraph{آماده‌سازی حالت پایه در شبیه‌ساز دیجیتال.}
%مسیر $H(s)$ را تروتر می‌کنند؛ گاف مسیر عمق مدار را تعیین می‌کند.
%سؤال: برآورد گاف روی خودِ دستگاه، با بودجه شات محدود، آیا زمان‌بندی را بهتر از حدس خطی می‌کند؟
%
%\paragraph{طیف‌سنجی مدل اسپینی روی کیوبیت‌های ابررسانا یا اتم خنثی.}
%خروجی آزمایش قله فوریه است.
%سؤال: کدام مشاهده‌پذیر محلی گاف هدف را نشان می‌دهد و کدام گاف تقارنی را پنهان می‌کند؟
%
%\paragraph{شیمی و مواد در فضای فعال کوچک.}
%گاف سینگلت--تریپلت یا شکافت دو حالت کم‌انرژی کمیت شیمیایی است~\cite{McArdle2020,GapHiggott2019}.
%سؤال: آیا \lr{VQD}/\lr{QSE} روی فضای فعال با گواهی خطا، در برابر \lr{CASCI} کلاسیک، فقط بازتولید است یا در مدلی که علامت مونت‌کارلو سخت است مزیت دارد؟
%
%\paragraph{حفاظت گاف در حافظه و گیت آدیاباتیک.}
%گاف بزرگ‌تر نرخ تحریک را کم می‌کند، ولی اجرای طولانی‌تر ناهمدوسی را زیاد می‌کند.
%سؤال کمّی: برای یک مدل نوفه‌دار مشخص، آیا $T$ بزرگ‌تر به‌خاطر گاف بهتر است یا $T$ کوتاه‌تر با \lr{CD} تقریبی؟
%
%\section{سختی مسئله و آنچه نباید ادعا کرد}
%
%تعیین انرژی پایه هامیلتونی محلی با دقت معکوس چندجمله‌ای در بدترین حالت \lr{QMA}-کامل است~\cite{GapGharibian2015}.
%تصمیم‌گیری گاف در حد ترمودینامیکی برای برخی خانواده‌های ساختگی حتی حل‌ناپذیر است~\cite{GapCubitt2015}.
%این دو نتیجه نمی‌گویند هر زنجیره آیزینگ سخت است، و مانع محاسبه دقیق ماتریس متناهی معلوم نمی‌شوند.
%در مقاله باید نمونه فیزیکی از بدترین حالت پیچیدگی جدا بماند.
%همچنین $n$ کیوبیت به‌تنهایی دشواری نیست: رژیم چندجمله‌ای در برابر نمایی را بخش~\ref{sec:complexity-scaling} جدا کرد.
%
%\section{مرور هدفمند منابع برای شروع پژوهش}
%
%این فهرست جایگزین بخش~\ref{sec:gap-papers} نیست؛ آن بخش جزئیات مقالات ۲۰۲۵--۲۰۲۶ را برای مسیر آدیاباتیک دارد.
%اینجا لنگرهای مروری و روش‌ساز آمده‌اند.
%\begin{itemize}
%	\item \textbf{شبیه‌سازی کوانتومی:} \lr{Georgescu--Ashhab--Nori} و چشم‌انداز مزیت عملی~\cite{Georgescu2014,Daley2022,CiracZoller2012}.
%	\item \textbf{آدیاباتیک:} \lr{Farhi et al.}، مرور \lr{Albash--Lidar}، کران \lr{Jansen--Ruskai--Seiler}، زمان‌بندی \lr{Roland--Cerf}~\cite{Farhi2000,AlbashLidar2018,GapJansen2007,GapRoland2002}.
%	\item \textbf{طیف کلاسیک:} مرور \lr{DMRG/MPS}~\cite{GapSchollwock2011,Cirac2021}.
%	\item \textbf{وردشی:} \lr{VQE}، مرور \lr{Tilly}، \lr{VQD}~\cite{Peruzzo2014,Tilly2022,GapHiggott2019}.
%	\item \textbf{زیرفضا و کریلوف:} \lr{QSE}، تحلیل کریلوف، اجرای پردازنده، نمونه‌برداری کریلوف~\cite{McClean2017QSE,Kirby2023Krylov,KrylovNatComm2025,GapYu2025}.
%	\item \textbf{میان‌بر آدیاباتیک:} \lr{Berry}، \lr{Sels--Polkovnikov}، نبرد با گاف نمایی، دیجیتال--آنالوگ~\cite{Berry2009STA,GapSels2017,Gueneau2025CD,Kumar2024DACQO}.
%	\item \textbf{برآورد گاف با سیگنال و \lr{FTQC}:}~\cite{GapLee2025,GapCugini2025,GapCarrera2026}.
%	\item \textbf{سختی:}~\cite{GapGharibian2015,GapCubitt2015}.
%\end{itemize}
%
%\section{مسئله‌های باز و طرح مقاله}
%\label{sec:gap-paper-tracks}
%
%یک مقاله در این حوزه وقتی شانس دارد که \emph{یک} ادعا، \emph{یک} خانواده مدل، و \emph{یک} معیار هزینه را هم‌زمان محدود کند.
%سه مسیر با فرضیه قابل‌آزمون، سازگار با فعالیت گروه (شبیه‌سازی، آیزینگ، آدیاباتیک، \lr{Qiskit}):
%
%\begin{enumerate}
%	\item \textbf{برآورد گاف مسیر با بودجه شات، سپس زمان‌بندی.}
%	فرضیه: روی زنجیره \lr{TFIM} با $n$ متوسط، تخمین نوسانی یا زیرفضاییِ $\Delta(s)$ با خطای نسبی معلوم، احتمال موفقیت را نسبت به جاروب خطی در زمان کل برابر بهبود می‌دهد.
%	مرجع: قطری‌سازی یا فرمیون آزاد برای اندازه‌های کوچک~\cite{GapPfeuty1970}؛ پروتکل:~\cite{GapLee2025,GapCugini2025}.
%	خروجی مقاله: منحنی $P_{\mathrm{succ}}(T)$ در برابر خطای گاف، نه فقط یک عدد $\Delta_{\min}$.
%	\item \textbf{تقریب موضعی رانش جبران‌کننده در برابر تغییر مسیر.}
%	فرضیه: برای همان خانواده، \lr{CD} محلی با شعاع محدود، نزدیک گاف چندجمله‌ای مفید است و نزدیک گاف نمایی اشباع می‌شود؛ کاتالیزگر ممکن است در رژیم دیگر بهتر باشد~\cite{GapSels2017,Gueneau2025CD}.
%	معیار: وفاداری نهایی در انرژی کنترل یا عمق مدار برابر.
%	\item \textbf{گواهی گاف در زیرفضای کریلوف، نه فقط تفاضل وردشی.}
%	فرضیه: با کنترل باقیمانده و عدد وضعیت $S$، بازه $[ \Delta_{\min}^{\mathrm{lo}}, \Delta_{\min}^{\mathrm{hi}} ]$ روی مدل اسپینی دوبعدی کوچک قابل‌گزارش است، در حالی که تفاضل دو \lr{VQE} بازه نمی‌دهد~\cite{Kirby2023Krylov,KrylovNatComm2025,McClean2017QSE}.
%	این مسیر به \lr{early-FTQC} نزدیک‌تر است و از صرفاً بازتولید \lr{VQD} فاصله می‌گیرد.
%\end{enumerate}
%
%مسیرهای~\ref{sec:gap-projects} را می‌توان زیرمجموعه‌ای از همین سه ادعا دید.
%دو اسپین فقط راستی‌آزمایی است و به‌تنهایی مقاله نمی‌سازد.
%
%\subsection{چک‌لیست قبل از شروع محاسبه}
%
%\begin{enumerate}
%	\item کدام گاف؟ کل، تقارنی، اپتیکی، مسیر $\Delta_{\min}(s)$؟
%	\item واحد انرژی چیست و با $n$ ثابت می‌ماند؟
%	\item مرجع کلاسیک تا چه $n$ دقیق است؟
%	\item هزینه را با چه واحدی می‌سنجیم: زمان فیزیکی، عمق تروتر، تعداد شات، یا پالس آنالوگ؟
%	\item نویز مدل شده است یا فقط سامانه بسته؟
%	\item ادعا بهبود $\Delta$ است یا بهبود $P_{\mathrm{succ}}$ در منابع برابر؟
%\end{enumerate}
%
%\section{ارتباط}
%
%تعریف گاف و نقش تقارن در بخش~\ref{sec:gap-research}؛ از گاف تا زمان در~\ref{sec:gap-runtime}؛ مقیاس‌بندی در~\ref{sec:complexity-scaling}؛ روش‌های عددی مسیر دو اسپین در~\ref{sec:gap-methods}--\ref{sec:gap-projects}.
%ابزار دیجیتال در فصل~\ref{chap:digital-tools}؛ \lr{QPE}/\lr{VQE}/\lr{QAOA} در فصل~\ref{chap:spectral-tools}؛ بستر آنالوگ در فصل~\ref{chap:hardware}.
%مدل آیزینگ در بخش~\ref{sec:ising}.
%
%\section{چه وقت استفاده می‌شود}
%
%وقتی سؤال پژوهش «گاف چیست؟»، «چطور آن را روی دستگاه بسنجیم؟»، یا «چطور با منابع محدود از گاف کوچک عبور کنیم؟» باشد.
%اگر فقط می‌خواهید یک مسیر آدیاباتیک مشخص را بفهمید، همان فصل~\ref{chap:evo-adi} کافی است؛ این فصل برای انتخاب خط مقاله است.
%
%\begin{remark}[جای توسعه]
%یک بنچمارک مشترک روی \lr{TFIM} باز و یک مدل با گاف نمایی، با همان بودجه شات؛ اتصال گاف لیوویلی به نرخ موفقیت در معادله مادر؛ مقایسه صریح \lr{QAOA} با مسیر آدیاباتیک از نظر گاف مؤثر، نه فقط مقدار نهایی هزینه.
%\end{remark}


% ============================================================
\part{گاف انرژی: برآورد، مهندسی، پژوهش}
% ============================================================

\chapter{گاف انرژی در شبیه‌سازی کوانتومی}
\label{chap:gap-energy}

\section{جایگاه در نقشه}

این فصل یک موضوع پژوهشی است، نه فقط یک کمیت کمکی برای فصل~\ref{chap:evo-adi}.
در نقشه شکل~\ref{fig:research-map} گاف در چند جایگاه متفاوت ظاهر می‌شود: به‌عنوان یک پرسش طیفی، به‌عنوان گلوگاه یک مسیر آدیاباتیک، و به‌عنوان یک ویژگی فیزیکی خودِ مدل.
بنابراین، اگرچه گاف در فصل~\ref{chap:evo-adi} عمدتاً از دیدگاه اثر آن بر اجرای آدیاباتیک بررسی شد، در این فصل خودِ گاف موضوع مطالعه است.

بخش‌های~\ref{sec:gap-research}--\ref{sec:gap-projects} گاف را در چارچوب یک مسیر $H(s)$ و مثال دوکیوبیتی باز می‌کنند.
اینجا همان کمیت را از آن مسیر جدا می‌کنیم و می‌پرسیم: در شبیه‌سازی کوانتومی چرا گاف مهم است، چگونه می‌توان آن را برآورد کرد، چگونه می‌توان رفتار آن را در یک مدل یا مسیر تغییر داد، و کدام مسئله‌های باز می‌توانند به یک مسئله پژوهشی قابل‌آزمون تبدیل شوند.

در این فصل، شبیه‌سازی آدیاباتیک یکی از کاربردهای مهم گاف در نظر گرفته می‌شود، نه چارچوب تعریف‌کننده آن.
پرسش اصلی، رفتار، برآورد و مهندسی گاف در مسئله‌های طیفی و دینامیکی شبیه‌سازی کوانتومی است؛ آدیاباتیک یکی از حالت‌های خاصی است که در آن گاف، علاوه بر معنای فیزیکی خود، مستقیماً بر هزینه دینامیکی آماده‌سازی اثر می‌گذارد.

هدف عملی فصل این است که پژوهش جاری گروه روی گاف متمرکز شود: مسئله، روش، معیار موفقیت، هزینه منابع، و رقیب کلاسیک از ابتدا روشن باشند.
از این دیدگاه، «گاف» به‌تنهایی موضوع مقاله نیست؛ سؤال پژوهشی باید مشخص کند که گاف را می‌خواهیم \emph{برآورد} کنیم، \emph{تفسیر} کنیم، یا \emph{مهندسی و بهره‌برداری} کنیم.

\section{ایده شهودی}

گاف انرژی به‌طور کلی فاصله میان دو تراز انرژی مورد نظر در طیف یک سامانه است.
در بسیاری از مسائل کم‌انرژی، مهم‌ترین تعریف فاصله میان حالت پایه و نخستین حالت برانگیخته است:

\begin{equation}
	\Delta = E_1-E_0.
	\label{eq:general-spectral-gap}
\end{equation}

با این حال، همین تعریف ساده در مسئله‌های مختلف می‌تواند به پرسش‌های متفاوتی منجر شود.
اگر هامیلتونی ثابت باشد، $\Delta$ یک ویژگی طیفی همان هامیلتونی است.
اگر هامیلتونی در امتداد یک مسیر تغییر کند، مثلاً $H(s)$، گاف نیز تابع مسیر می‌شود:

\begin{equation}
	\Delta(s)=E_1(s)-E_0(s),
	\qquad s\in[0,1].
	\label{eq:instantaneous-gap-general}
\end{equation}

در این حالت کمینه گاف روی مسیر برابر است با

\begin{equation}
	\Delta_{\min}
	=
	\min_{s\in[0,1]}\Delta(s).
	\label{eq:minimum-gap-general}
\end{equation}

این سه کمیت را باید از یکدیگر جدا نگه داشت:

\begin{itemize}
	\item $\Delta$: گاف یک هامیلتونی در یک نقطه مشخص از فضای پارامترها؛
	\item $\Delta(s)$: گاف لحظه‌ای در امتداد یک مسیر؛
	\item $\Delta_{\min}$: کوچک‌ترین گاف در کل مسیر.
\end{itemize}

این تمایز در فصل حاضر اهمیت اساسی دارد، زیرا $\Delta$ می‌تواند یک ویژگی فیزیکی مدل باشد، در حالی که $\Delta_{\min}$ در یک مسئله آدیاباتیک به یک ویژگی مشترک مدل، مسیر و انتخاب پارامترسازی تبدیل می‌شود.

گاف عدد کوچکی به نظر می‌رسد، ولی نقش‌های کاملاً متفاوتی بازی می‌کند:

\begin{itemize}
	\item در محاسبه آدیاباتیک، گاف کمینه مسیر یکی از عوامل اصلی تعیین‌کننده زمان تحول است (بخش~\ref{sec:gap-runtime}).
	\item در فیزیک بس‌ذره‌ای، بسته یا باز بودن گاف می‌تواند نشانه‌ای از ساختار فازی سامانه باشد.
	\item در طیف‌سنجی، گاف یک مقیاس انرژی یا فرکانس مشخص می‌کند که باید از داده زمانی یا پاسخ سامانه استخراج شود.
	\item در برآورد انرژی‌های برانگیخته، دقت تخمین دو ویژه‌مقدار مستقیماً بر دقت گاف اثر می‌گذارد.
	\item در برخی سامانه‌های باز، گاف هامیلتونی بسته تنها مقیاس مرتبط با دینامیک نیست و باید گاف مولد دینامیک باز را از آن جدا کرد.
\end{itemize}

بنابراین دو مسئله متفاوت باید از ابتدا از یکدیگر جدا شوند:

\begin{enumerate}
	\item \textbf{برآورد گاف:} گاف یک مدل مشخص چیست و چگونه می‌توان آن را با منابع محدود اندازه گرفت یا محاسبه کرد؟
	\item \textbf{مهندسی گاف و دینامیک:} آیا می‌توان با تغییر مدل، مسیر، رمزگذاری یا پروتکل کنترل، رفتار گاف یا اثر آن بر دینامیک را بهبود داد؟
\end{enumerate}

اولی یک \emph{پرسش طیفی} است و دومی یک \emph{مسئله مهندسی} است.
شکل~\ref{fig:gap-methods-map} این تفکیک را نشان می‌دهد.

\begin{figure}[htbp]
	\centering
	\begin{latin}
		\setstretch{1}
		\resizebox{\ifdim\width>\textwidth\textwidth\else\width\fi}{!}{%
			% دسته‌بندی روش‌های برآورد و مهندسی گاف
% سبک مشترک نمودارهای سند پایه‌ها (فقط داخل محیط latin فراخوانی شود)
\tikzset{
	qafbox/.style={
		draw=black!55,
		rounded corners=2pt,
		align=center,
		inner sep=3pt,
		thick,
		fill=white,
		font=\footnotesize\linespread{1}\selectfont,
		minimum height=0.95cm
	},
	qafwide/.style={
		qafbox,
		text width=2.55cm
	},
	qafnarr/.style={
		qafbox,
		text width=2.05cm
	},
	qaftitle/.style={
		font=\small\bfseries\linespread{1}\selectfont,
		align=center
	},
	qafarr/.style={
		-{Latex[length=2.2mm]},
		thick,
		draw=black!65
	},
	qafpanel/.style={
		draw,
		thick,
		rounded corners=4pt,
		inner xsep=8pt,
		inner ysep=8pt
	}
}

\begin{tikzpicture}[
	x=1cm,
	y=1cm,
	font=\footnotesize,
	gbox/.style={qafbox, text width=2.35cm, minimum height=1.05cm}
	]
	\node[gbox, fill=blue!12, text width=3.4cm] (why) at (0,3.4) {Why the gap?};
	\node[gbox, fill=green!12] (cl) at (-5.2,1.6) {Classical spectrum};
	\node[gbox, fill=orange!14] (dyn) at (-1.75,1.6) {Dynamics / spectroscopy};
	\node[gbox, fill=cyan!14] (var) at (1.75,1.6) {Variational / subspace};
	\node[gbox, fill=red!12] (ft) at (5.2,1.6) {FTQC / QPE};
	\node[gbox, fill=green!10, text width=2.55cm] (cl2) at (-5.2,0) {Lanczos, DMRG, QMC};
	\node[gbox, fill=orange!10, text width=2.55cm] (dyn2) at (-1.75,0) {Oscillation, quench, $C(\tau)$};
	\node[gbox, fill=cyan!10, text width=2.55cm] (var2) at (1.75,0) {VQD, QSE, Krylov};
	\node[gbox, fill=red!10, text width=2.55cm] (ft2) at (5.2,0) {Phase est., counting};
	\node[gbox, fill=violet!12, text width=4.6cm] (eng) at (0,-1.7) {Gap engineering / shortcuts};
	\node[font=\scriptsize, text=black!60] at (0,-2.45) {schedule, catalyst, CD, encoding};
	\draw[qafarr] (why) -- (cl);
	\draw[qafarr] (why) -- (dyn);
	\draw[qafarr] (why) -- (var);
	\draw[qafarr] (why) -- (ft);
	\draw[qafarr] (cl) -- (cl2);
	\draw[qafarr] (dyn) -- (dyn2);
	\draw[qafarr] (var) -- (var2);
	\draw[qafarr] (ft) -- (ft2);
	\draw[qafarr] (cl2.south) -- ++(0,-0.35) -| (eng.north);
	\draw[qafarr] (dyn2.south) -- ++(0,-0.35) -| (eng.north);
	\draw[qafarr] (var2.south) -- ++(0,-0.35) -| (eng.north);
	\draw[qafarr] (ft2.south) -- ++(0,-0.35) -| (eng.north);
\end{tikzpicture}
%
		}
	\end{latin}
	\caption{چهار خانواده اصلی برای برآورد گاف و یک خانواده روش‌های مهندسی. برآورد می‌پرسد گاف چیست و چگونه می‌توان آن را استخراج کرد؛ مهندسی می‌پرسد چگونه می‌توان مسیر، هامیلتونی یا کنترل را برای دستیابی به عملکرد بهتر تغییر داد.}
	\label{fig:gap-methods-map}
\end{figure}

\section{چرا گاف در این حوزه مهم است؟}

\subsection{چند معنای فیزیکی که نباید قاطی شوند}

کلمه «گاف» در مقالات مختلف الزاماً به یک کمیت واحد اشاره نمی‌کند.
برای یک پژوهش قابل‌انتشار باید از ابتدا مشخص شود کدام گاف، بین کدام ترازها، در کدام زیرفضای فیزیکی، و با کدام مشاهده‌پذیر یا فرآیند دینامیکی مد نظر است~\cite{Georgescu2014,AlbashLidar2018}.

در ادامه چند معنای مهم از گاف از یکدیگر جدا می‌شوند:

\begin{enumerate}
	\item \textbf{گاف طیفی کم‌انرژی.}
	برای یک هامیلتونی بسته، گاف میان حالت پایه و نخستین حالت برانگیخته به‌صورت
	$\Delta=E_1-E_0$
	تعریف می‌شود.
	اگر تقارن یا عدد کوانتومی خاصی حفظ شود، منظور از $E_1$ باید نخستین حالت برانگیخته \emph{قابل‌دسترسی در همان زیرفضا} باشد.
	همین مفهوم در فصل~\ref{chap:evo-adi} برای مسیر آدیاباتیک به‌کار رفت.
	

	\item \textbf{گاف بخش تقارنی.}
	اگر
	$[H,S]=0$
	باشد، فضای هیلبرت به بخش‌های تقارنی تقسیم می‌شود.
	در این شرایط ممکن است دو تراز متعلق به دو بخش متفاوت به یکدیگر برسند، بدون آنکه گذار میان آنها تحت دینامیک مورد نظر مجاز باشد.
	در نتیجه، گاف کل می‌تواند صفر شود، در حالی که گاف مربوط به بخش قابل‌دسترسی مثبت باقی بماند (بخش~\ref{sec:gap-research}).
	
	\item \textbf{گاف بار، اسپینی یا اپتیکی.}
	در مدل‌هایی مانند هابارد یا در مسائل مولکولی، فاصله دو ویژه‌مقدار کم‌انرژی الزاماً همان انرژی لازم برای افزودن یک الکترون، ایجاد یک برانگیختگی اسپینی، یا جذب یک فوتون نیست.
	قواعد انتخاب و عنصر ماتریسی مشاهده‌پذیر تعیین می‌کنند کدام شکافت در یک آزمایش یا پروتکل اندازه‌گیری ظاهر شود.
	
	\item \textbf{گاف توپولوژیک.}
	در سامانه‌هایی که فاز توپولوژیک دارند، گاف حجمی بخشی از ساختار فازی سامانه است و بسته‌شدن آن می‌تواند با تغییر فاز همراه شود.
	در عین حال، حالت‌های لبه یا مرزی ممکن است گاف کوچک یا صفر داشته باشند و بنابراین گاف حجمی را نباید با همه شکافت‌های طیف یکی دانست.
	
	\item \textbf{گاف لیوویلی.}
	در سامانه باز، نرخ‌های واهلش و زمان‌های مشخصه دینامیک را باید از طیف مولد دینامیک باز استخراج کرد.
	در این حالت، گاف مربوط به مولد لیوویلی الزاماً برابر با $E_1-E_0$ هامیلتونی بسته نیست~\cite{GapSels2017}.
\end{enumerate}
بنابراین، در فصل حاضر کمیت مرکزی همان گاف طیفی هامیلتونی است و سایر تعاریف برای مشخص‌کردن مرز کاربردها یا نشان‌دادن تفاوت میان مسئله‌های فیزیکی وارد می‌شوند.
این تفکیک برای طراحی یک مسئله ژوهشی ضروری است؛ اگر مقاله بدون مشخص‌کردن اینکه کدام گاف را اندازه می‌گیرد یا مهندسی می‌کند از واژه «gap» استفاده کند، نتیجه برای مقایسه و داوری قابل تفسیر نخواهد بود.

\subsection{از گاف یک هامیلتونی تا گاف یک مسیر}

در یک هامیلتونی ثابت، پرسش طیفی ساده است:
\[
\Delta=E_1-E_0.
\]
اما در یک مسئله دینامیکی یا آدیاباتیک، هامیلتونی می‌تواند به یک پارامتر کنترل وابسته باشد:
\[
H=H(s).
\]
در این حالت طیف نیز با $s$ تغییر می‌کند:
\begin{equation}
	E_j=E_j(s),
	\qquad
	\Delta(s)=E_1(s)-E_0(s).
\end{equation}
بنابراین $\Delta_{\min}$ دیگر تنها ویژگی یکی از هامیلتونی‌های ابتدا یا انتهای مسیر نیست، بلکه ویژگی کل مسیر در فضای هامیلتونی‌هاست:
\begin{equation}
	\Delta_{\min}
	=
	\min_{s\in[0,1]}
	\left[E_1(s)-E_0(s)\right].
\end{equation}
این تمایز، ارتباط فصل حاضر با فصل~\ref{chap:evo-adi} را روشن می‌کند.
در آن فصل سؤال اصلی این بود که چگونه یک مسیر مشخص را با احتمال موفقیت مناسب طی کنیم؛ در این فصل همان گاف از دیدگاه وسیع‌تر بررسی می‌شود: چگونه آن را پیدا کنیم، چگونه بفهمیم کمینه آن از کجا ناشی می‌شود، و آیا می‌توان با تغییر مسیر یا روش اندازه‌گیری، رفتار آن را بهبود داد.

\subsection{نقش گاف در محاسبه آدیاباتیک}
\label{sec:gap-adiabatic-role}

در رهیافت آدیاباتیک، گاف کمینه یکی از مهم‌ترین عوامل تعیین‌کننده هزینه تحول زمانی است~\cite{Farhi2000,GapJansen2007,AlbashLidar2018}.
به‌صورت مقیاسی، در بسیاری از شرایط می‌توان وابستگی اصلی را به شکل
\begin{equation}
	T\propto \frac{1}{\Delta_{\min}^{2}}
	\label{eq:gap-runtime-scaling}
\end{equation}
بیان کرد، مشروط بر اینکه مقیاس انرژی هامیلتونی و سایر عوامل دینامیکی در مقایسه ثابت نگه داشته شوند.
این رابطه یک قانون عمومی با ضریب واحد نیست؛ در کران‌های دقیق، مشتق‌های هامیلتونی و عناصر ماتریسی گذار نیز وارد می‌شوند.
بنابراین در مقایسه دو پروتکل باید علاوه بر $\Delta_{\min}$، مقیاس انرژی، مسیر و کنترل دینامیکی نیز مشخص باشند.

بسته‌شدن چندجمله‌ای در برابر نمایی همان تمایزی است که بخش~\ref{sec:complexity-scaling} باز کرد.
اگر، برای یک خانواده مدل با مقیاس انرژی ثابت،
\begin{equation}
	\Delta_{\min}(n)\sim n^{-\alpha},
\end{equation}
آنگاه رابطه  مقیاسی~\eqref{eq:gap-runtime-scaling} به یک هزینه چندجمله‌ای منجر می‌شود.
در مقابل، اگر
\begin{equation}
	\Delta_{\min}(n)\sim e^{-cn},
\end{equation}
هزینه زمانی متناظر می‌تواند به‌صورت نمایی رشد کند.

پژوهش گاف در این شاخه سه سؤال متمایز دارد:
\begin{itemize}
	\item گاف لحظه‌ای $\Delta(s)$ در یک مدل مشخص چیست؟
	\item کمینه آن روی مسیر کجاست و با $n$ چگونه بسته می‌شود؟
	\item آیا با عوض کردن مسیر، زمان‌بندی، رمزگذاری، یا جمله کمکی می‌توان اثر گاف کوچک را کاهش داد؟
\end{itemize}

سؤال سوم الزاماً به معنای «بزرگ‌کردن گاف» نیست.
ممکن است یک روش، بدون تغییر قابل‌توجه در $\Delta_{\min}$، تحریک‌های دیاباتیک را سرکوب کند.
این تمایز در بخش~\ref{sec:gap-eng} دنبال می‌شود.

\subsection{نقش گاف در شبیه‌سازی دینامیکی و مواد}

در شبیه‌سازی تحول زمانی، گاف می‌تواند به‌صورت یک مقیاس فرکانسی در سیگنال‌های زمانی ظاهر شود.
اگر یک حالت اولیه با حالت پایه و یک حالت برانگیخته هم‌پوشانی داشته باشد و مشاهده‌پذیر $O$ بین آن دو عنصر ماتریسی غیرصفر داشته باشد، جمله‌ای از مقدار مورد انتظار آن مشاهده‌پذیر به شکل

\begin{equation}
	\langle O(t)\rangle
	=
	A
	+
	2\mathrm{Re}
	\left[
	c_0^*c_1O_{01}
	e^{-i\Delta t/\hbar}
	\right]
	+\cdots
\end{equation}

ظاهر می‌شود.
بنابراین گاف انرژی به یک فرکانس تبدیل می‌شود:

\begin{equation}
	\omega=\frac{\Delta}{\hbar}.
\end{equation}

این رابطه یک مسیر مهم برای برآورد گاف از روی داده دینامیکی فراهم می‌کند.
از دید عملی نیز تفکیک دو فرکانس نزدیک نیازمند مشاهده زمانی کافی است؛ به‌طور تقریبی،

\begin{equation}
	\delta E\sim \frac{2\pi\hbar}{t_{\max}},
\end{equation}

به این معنا که گاف کوچک‌تر می‌تواند به بازه زمانی طولانی‌تر و در نتیجه به کنترل همدوسی بهتر نیاز داشته باشد~\cite{Georgescu2014,Daley2022}.

در شیمی کوانتومی، گاف هومو-لومو یا شکافت سینگلت-تریپلت می‌تواند یک کمیت فیزیکی مستقل باشد، نه صرفاً گلوگاه یک الگوریتم~\cite{McArdle2020,Bauer2020}.
در شبیه‌ساز آنالوگِ شبکه نوری یا اتم خنثی نیز طیف‌سنجی پس از \lr{quench} یکی از راه‌های مطالعه طیف است؛ در این حالت گاف می‌تواند از مکان قله‌های ساختار دینامیکی استخراج شود، بدون آنکه لازم باشد کل ماتریس هامیلتونی به‌صورت کلاسیک قطری‌سازی شود.

این مثال‌ها نشان می‌دهند که گاف در شبیه‌سازی کوانتومی تنها وسیله‌ای برای تخمین زمان اجرای یک الگوریتم نیست؛ خودِ گاف می‌تواند بخشی از پاسخ فیزیکی مورد نظر باشد.

\subsection{نقش گاف در کیوبیت و سخت‌افزار}

برای کیوبیت ابررسانا، فاصله انرژی میان ترازهای پایین و ترازهای بالاتر بخشی از طراحی فیزیکی سامانه را تعیین می‌کند (فصل~\ref{chap:models}).
گاف بزرگ‌تر می‌تواند در برابر برخی تحریک‌های گرمایی محافظت بیشتری ایجاد کند، مشروط بر آنکه عملگر جفت‌شدگی محیط و طیف نویز واقعاً آن گذار را تحریک کنند.

در سامانه‌های حفاظت‌شده و کدهای پایدارکننده نیز می‌توان از هامیلتونی‌های جریمه برای ایجاد شکافت میان زیرفضای محاسباتی و حالت‌های خطادار استفاده کرد.
در اینجا نیز گاف به‌تنهایی تعیین‌کننده عملکرد نیست و باید آن را همراه با ساختار خطا، نحوه کوپل محیط و هزینه پیاده‌سازی بررسی کرد.

در نتیجه، پژوهش گاف فقط «الگوریتم آدیاباتیک» نیست.
طراحی کیوبیت، حافظه، شبیه‌ساز آنالوگ، طیف‌سنجی و آماده‌سازی حالت همگی می‌توانند از برآورد یا مهندسی گاف بهره ببرند.

\section{رهیافت‌های اساسی برای یافتن گاف}

برای تبدیل گاف به یک موضوع پژوهشی عملی، نخست باید مشخص شود که آن را چگونه می‌توان برآورد کرد.
چهار خانواده اصلی در این فصل از یکدیگر جدا می‌شوند:

\begin{enumerate}
	\item محاسبه مستقیم طیف با روش‌های کلاسیک؛
	\item طیف‌سنجی دینامیکی از روی داده زمانی؛
	\item روش‌های وردشی و زیرفضایی روی سخت‌افزار کوانتومی؛
	\item برآورد فاز و روش‌های تحمل‌پذیر در برابر خطا.
\end{enumerate}

جدول~\ref{tab:gap-method-families} این خانواده‌ها را با رژیم معمول و ایده اصلی آنها جمع می‌کند.
جزئیات عددی مسیر دوکیوبیتی در بخش~\ref{sec:gap-methods} آمده است؛ در اینجا هر خانواده به‌عنوان یک خط پژوهشی بررسی می‌شود.

\begin{table}[htbp]
	\centering
	\small
	\setstretch{1.12}
	\caption{چهار خانواده اصلی برآورد گاف در شبیه‌سازی کوانتومی. هزینه‌ها نوعی‌اند و کران قضیه محسوب نمی‌شوند.}
	\label{tab:gap-method-families}
	\begin{tabularx}{\textwidth}{@{}
			>{\raggedright\arraybackslash}p{3.3cm}
			>{\raggedright\arraybackslash}X
			>{\raggedright\arraybackslash}p{3.4cm}
			@{}}
		\toprule
		\rowcolor{gray!14}
		\textbf{خانواده} & \textbf{ایده} & \textbf{رژیم نوعی} \
		\midrule
		طیف کلاسیک
		& قطری‌سازی، لانچوس، \lr{DMRG}، \lr{QMC}
		& $n$ کوچک تا متوسط یا درهم‌تنیدگی محدود
		\
		\rowcolor{gray!8}
		طیف‌سنجی دینامیکی
		& نوسان $\langle O(t)\rangle$، \lr{quench}، همبستگی زمان موهومی
		& آنالوگ و دیجیتال؛ وابسته به مشاهده‌پذیر
		\
		وردشی / زیرفضا
		& \lr{VQD}، \lr{SSVQE}، \lr{QSE}، کریلوف کوانتومی
		& \lr{NISQ} و \lr{early-FTQC}
		\
		\rowcolor{gray!8}
		برآورد فاز / شمارش
		& \lr{QPE}، \lr{QSVT}، پرسمان شمارش
		& عمدتاً در افق \lr{FTQC}
		\
		\bottomrule
	\end{tabularx}
\end{table}

\subsection{طیف کلاسیک: مرجع، نه رقیب همیشه مغلوب}

برای هر مقاله روش کوانتومی، مرجع کلاسیک باید صریح باشد~\cite{GapSchollwock2011,Cirac2021}.
قطری‌سازی کامل برای $D=2^n$ با حافظه $O(D^2)$ برای اندازه‌های کوچک مناسب است، اما هزینه آن با افزایش $n$ به‌سرعت رشد می‌کند.
روش لانچوس با استفاده از ضرب $Hv$ بدون ساختن ماتریس متراکم می‌تواند چند تراز پایین را استخراج کند؛ همگرایی آن به گاف، بردار شروع و ساختار طیف وابسته است.
شبکه تانسوری (\lr{DMRG/MPS}) برای زنجیره‌های کم‌بعد و کم‌درهم‌تنیده همچنان یک مرجع کلاسیک قدرتمند است.
مونت‌کارلو کوانتومی نیز در مدل‌هایی که با مشکل علامت مواجه نیستند می‌تواند از رفتار همبستگی در زمان موهومی اطلاعات طیفی استخراج کند.

\textbf{نتیجه برای پروژه.}
اگر مدل ما زنجیره \lr{TFIM} یا یک مولکول کوچک باشد، مرجع کلاسیک غالباً دقیق‌تر و ارزان‌تر از مدار نوفه‌دار کوانتومی است.
در چنین شرایطی، صرف بازتولید یک عدد گاف با سخت‌افزار کوانتومی به‌تنهایی مزیت محاسباتی ایجاد نمی‌کند.
سهم کوانتومی زمانی معنادارتر است که مدل در رژیم سخت کلاسیک قرار گیرد، یا خودِ \emph{پروتکل اندازه‌گیری و برآورد گاف روی سخت‌افزار} موضوع پژوهش باشد.

بنابراین پرسش «چرا کوانتومی؟» باید در کنار پرسش «گاف چیست؟» از ابتدا پاسخ داده شود.

\subsection{طیف‌سنجی دینامیکی: گاف به‌عنوان فرکانس}

اگر حالت آزمایشی هم‌پوشانی با حالت پایه و یک حالت برانگیخته داشته باشد و مشاهده‌پذیر عنصر ماتریسی غیرصفر داشته باشد،

\begin{equation}
	\langle O(t)\rangle
	=
	A
	+
	2,\mathrm{Re}\bigl[
	c_0^*c_1O_{01}e^{-i\Delta t/\hbar}
	\bigr]
	+\cdots
	\label{eq:gap-osc-research}
\end{equation}

فرکانس نوسان، در صورت غالب‌بودن همین گذار، مستقیماً با گاف مرتبط است (رابطه~\eqref{eq:gap-oscillation} در بخش~\ref{sec:gap-methods}).
در شبیه‌ساز آنالوگ همین ایده به شکل \lr{quench} و تبدیل فوریه تابع ساختار دینامیکی درمی‌آید~\cite{Georgescu2014,Daley2022}.
در زمان موهومی نیز شیب $\log C_O(\tau)$ می‌تواند به کمترین گاف دارای وزن در مشاهده‌پذیر میل کند.

محدودیت اصلی این خانواده انتخاب $O$ است: نبود قله دلیل نبود تراز نیست.
تقارن می‌تواند یک گذار را از دید مشاهده‌پذیر خاموش کند و نویز خوانش می‌تواند قله را جابه‌جا یا پهن کند.
در نتیجه، یک پروتکل طیف‌سنجی باید علاوه بر استخراج فرکانس، \emph{انتخاب مشاهده‌پذیر، هم‌پوشانی حالت اولیه، وضوح زمانی، و مدل نویز} را نیز مشخص کند.
مقالات \lr{Lee--Myers--Scarola} و \lr{Cugini et al.} همین خط را با پردازش سیگنال و آماده‌سازی آدیاباتیک دنبال می‌کنند~\cite{GapLee2025,GapCugini2025}.

\subsection{وردشی و زیرفضا: دو انرژی، یک تفاضل خطرناک}

\lr{VQE} به‌طور مستقیم یک کران بالا برای $E_0$ فراهم می‌کند، نه گاف~\cite{Peruzzo2014,Tilly2022}.
برای استخراج $E_1$ باید قید یا سازوکار اضافی برای جداکردن حالت‌های برانگیخته از حالت پایه وارد شود.

در \lr{VQD}، جریمه هم‌پوشانی با حالت‌های پایین‌تر برای هدف‌گرفتن تراز بعدی به‌کار می‌رود~\cite{GapHiggott2019}:

\begin{equation}
	{\cal L}*k(\theta)
	=
	\langle\psi(\theta)|H|\psi(\theta)\rangle
	+
	\sum*{j<k}\mu_j
	\bigl|\langle\psi_j|\psi(\theta)\rangle\bigr|^2.
\end{equation}

\lr{SSVQE} چند حالت را هم‌زمان در یک زیرفضای پارامتری بهینه می‌کند.
بسط زیرفضای کوانتومی (\lr{QSE}) از تحریک‌های موضعی یک حالت مرجع، ماتریس‌های $A$ و $S$ می‌سازد و مسئله ویژه تعمیم‌یافته را کلاسیک حل می‌کند~\cite{McClean2017QSE}.
کریلوف کوانتومی نیز با استفاده از توان‌های هامیلتونی یا تحول زمانی یک حالت مرجع، یک زیرفضا برای استخراج اطلاعات طیفی می‌سازد؛ اجرای سخت‌افزاری این ایده در اندازه‌های قابل‌توجه نیز گزارش شده است~\cite{KrylovNatComm2025,GapYu2025,Kirby2023Krylov}.

\textbf{هشدار روش‌شناختی.}
تفاضل دو انرژی وردشی لزوماً یک کران تضمینی برای $\Delta$ نیست.
اگر خطای دو تخمین هم‌علامت یا ناهم‌علامت باشد، خطای گاف می‌تواند از خطای هر یک از انرژی‌ها کوچک‌تر یا بزرگ‌تر شود.
بنابراین هر ادعای مربوط به گاف در این خانواده باید، تا حد امکان، باقیمانده

\begin{equation}
	|H|\psi\rangle-E|\psi\rangle|
\end{equation}

همگرایی زیرفضا و در صورت امکان یک گواهی فاصله تا بخش مورد نظر طیف همراه شود.

\subsection{برآورد فاز و الگوریتم‌های تحمل‌پذیر}

اگر بتوان

\begin{equation}
	U=e^{-iH\tau/\hbar}
\end{equation}

را با دقت کنترل‌شده ساخت، \lr{QPE} ویژه‌مقدارهای $H$ را به فازهای قابل‌اندازه‌گیری تبدیل می‌کند~\cite{NielsenChuang2010,McArdle2020}.
اگر خطای مطلق دو انرژی به‌ترتیب $\epsilon_0$ و $\epsilon_1$ باشد، آنگاه خطای گاف از رابطه

\begin{equation}
	|\delta\Delta|
	\leq
	\epsilon_0+\epsilon_1
\end{equation}

تبعیت می‌کند و در حالت متقارن $\epsilon_0=\epsilon_1=\epsilon_E$ داریم

\begin{equation}
	|\delta\Delta|\leq2\epsilon_E.
\end{equation}

در اینجا زمان همدوس لازم برای تفکیک انرژی با دقت مشخص، تنها بخشی از هزینه است؛ هزینه ساخت و اجرای تکامل $U$ نیز باید جداگانه حساب شود.
الگوریتم‌های مبتنی بر بلوک‌کدگذاری و تبدیل تکین (\lr{QSVT}) نیز چارچوب‌هایی برای برآورد ویژگی‌های طیفی با تضمین پیچیدگی فراهم می‌کنند؛ اما هزینه بارگذاری داده، شبیه‌سازی هامیلتونی و دسترسی به منابع لازم باید جداگانه بررسی شود~\cite{GapCarrera2026}.

این خط پژوهشی عمدتاً به افق \lr{FTQC} مربوط است، نه لزوماً دستگاه نوفه‌دار ده‌کیوبیتی.
در نتیجه، مقایسه آن با روش‌های \lr{NISQ} باید بر اساس یک معیار مشترک انجام شود و نباید صرفاً بر اساس نام الگوریتم صورت گیرد.

\section{مهندسی گاف و مهندسی دینامیک}
\label{sec:gap-eng}

«گاف را بزرگ کنیم» با «گذار را کم کنیم» یکی نیست.
در یک مسئله آدیاباتیک یا دینامیکی، می‌توان به‌جای افزایش مستقیم $\Delta_{\min}$، مسیر، زمان‌بندی یا کنترل را تغییر داد تا عملکرد نهایی بهتر شود~\cite{AlbashLidar2018,GapSels2017,Gueneau2025CD}.

از این دیدگاه، چهار مداخله اصلی را می‌توان از یکدیگر جدا کرد:

\begin{enumerate}
	\item تغییر زمان‌بندی روی همان مسیر؛
	\item تغییر مسیر یا هامیلتونی؛
	\item افزودن کنترل‌های جبران‌کننده برای کاهش گذارهای ناخواسته؛
	\item تغییر رمزگذاری یا انتخاب زیرفضای فیزیکی.
\end{enumerate}

معیار موفقیت نیز نباید صرفاً «بزرگ‌ترشدن گاف» باشد.
اگر یک روش $\Delta_{\min}$ را افزایش دهد اما عمق مدار، تعداد گیت، مدت اجرای فیزیکی یا حساسیت به نویز را بسیار زیاد کند، ممکن است در عمل روش بهتری نباشد.

\subsection{زمان‌بندی: گاف را عوض نمی‌کند}

زمان‌بندی غیرخطی $s(t)$ همان مسیر $H(s)$ را با سرعت متغیر طی می‌کند.
در نتیجه، خود تابع طیفی

\begin{equation}
	\Delta(s)
\end{equation}

تغییر نمی‌کند؛ فقط مدت زمانی که سیستم در نواحی مختلف مسیر می‌گذراند تغییر می‌کند.
در نتیجه می‌توان نزدیک نواحی با گاف کوچک آهسته‌تر حرکت کرد و در نواحی کم‌حساسیت سریع‌تر پیش رفت~\cite{GapRoland2002}.

این روش از نظر پیاده‌سازی یکی از ساده‌ترین مداخلات است و باید خط پایه هر مقاله مهندسی باشد.
در مقایسه دو زمان‌بندی باید مسیر، زمان کل و منابع کنترل ثابت یا دقیقاً گزارش شوند.

\subsection{تغییر مسیر و هامیلتونی کاتالیزگر}

در این روش، خودِ مسیر در فضای هامیلتونی‌ها تغییر می‌کند.
برای مثال می‌توان جمله‌ای از شکل

\begin{equation}
	H(s)
	=
	H_0(s)
	+
	f(s)H_{\mathrm{cat}},
	\qquad
	f(0)=f(1)=0
\end{equation}

به مسیر اضافه کرد.

در این حالت برخلاف تغییر زمان‌بندی، خودِ طیف $\Delta(s)$ نیز می‌تواند تغییر کند.
هدف معمولاً اجتناب از نواحی با شکافت بسیار کوچک، حذف یا جابه‌جایی یک گذار بحرانی، یا افزایش $\Delta_{\min}$ است.
با این حال، بهای روش ممکن است افزایش غیرمحلی‌شدن هامیلتونی، افزایش تعداد جملات قابل‌کنترل یا افزایش پیچیدگی پیاده‌سازی باشد.

بنابراین بزرگ‌شدن گاف در یک نقطه از مسیر کافی نیست.
باید کمینه کل مسیر، عناصر ماتریسی گذار، و هزینه اجرای هامیلتونی جدید دوباره محاسبه شوند (بخش~\ref{sec:gap-control}).

\subsection{رانش جبران‌کننده: میان‌بر، نه لزوماً گاف بزرگ‌تر}

رانش جبران‌کننده\ft{counterdiabatic driving} یک جمله کمکی به دینامیک اضافه می‌کند تا گذارهای ناخواسته را سرکوب کند، حتی اگر گاف مسیر مرجع کوچک باقی بماند~\cite{Berry2009STA,GapSels2017}.
بنابراین این روش را نباید صرفاً یک روش «افزایش گاف» دانست.

جمله دقیق رانش جبران‌کننده معمولاً غیرمحلی است و به اطلاعاتی از ساختار ویژه‌مقدارها و ویژه‌بردارهای هامیلتونی وابسته است.
از این رو، استفاده مستقیم از آن می‌تواند بخشی از دشواری اصلی را به مسئله ساخت کنترل مناسب منتقل کند.
تقریب‌های موضعی و بسط وردشی پتانسیل پیمانه‌ای آدیاباتیک عملی‌ترند، ولی نزدیک گاف‌های بسیار کوچک باید اثر خطای تقریب و هزینه پیاده‌سازی به‌طور صریح سنجیده شود~\cite{Gueneau2025CD}.

پیاده‌سازی دیجیتال، دیجیتال--آنالوگ روی یون، و آنالوگ روی اتم خنثی نیز برای این ایده گزارش شده است~\cite{Kumar2024DACQO,ACQC2026}.

برای گروه ما سؤال پژوهشی دقیق می‌تواند این باشد:
در یک مدل اسپینی محلی با منابع کنترل محدود، تقریب موضعی \lr{CD} تا چه اندازه می‌تواند اثر یک $\Delta_{\min}$ کوچک را جبران کند، و هزینه گیت یا پالس آن در برابر اجرای یک جاروب آهسته‌تر چگونه مقایسه می‌شود؟

این سؤال عمداً «گاف» و «عملکرد دینامیکی» را هم‌زمان در نظر می‌گیرد.

\subsection{رمزگذاری و تقارن}

گاهی یک گاف کوچک نتیجه انتخاب نامناسب نمایش مسئله است، نه الزاماً ویژگی اجتناب‌ناپذیر فیزیک مورد نظر.
انتخاب هامیلتونی آغازین، حفظ یا شکستن تقارن، نگاشت مسئله به کیوبیت‌ها و انتخاب فضای محاسباتی می‌تواند طیف و در نتیجه $\Delta_{\min}(n)$ را تغییر دهد.

این شاخه به بهینه‌سازی ترکیبی نزدیک است، ولی اگر مدل فیزیکی اسپینی باشد همچنان می‌توان آن را در چارچوب شبیه‌سازی کوانتومی مطالعه کرد.
پرسش اصلی این است که آیا می‌توان همان مسئله فیزیکی را در نمایشی با رفتار طیفی مطلوب‌تر و هزینه پیاده‌سازی کمتر بازنمایی کرد.

در اینجا نیز باید بین «گاف فیزیکی مدل» و «گافی که در نتیجه یک رمزگذاری خاص ایجاد شده است» تفاوت گذاشت.

\section{کاربرد در سامانه‌های کیوبیتی و شبیه‌سازی}

گستره روش‌های برآورد و مهندسی گاف را می‌توان به چند سناریوی مشخص تبدیل کرد که هر یک ظرفیت یک مسئله پژوهشی مستقل را دارند.

\paragraph{آماده‌سازی حالت پایه در شبیه‌ساز دیجیتال.}
مسیر $H(s)$ را می‌توان به گام‌های زمانی تقسیم و با روش‌هایی مانند تروترایزیشن روی مدار دیجیتال پیاده کرد.
در اینجا گاف مسیر بر حساسیت فرآیند به زمان‌بندی اثر می‌گذارد و تقریب تروتر نیز یک منبع خطای مستقل است.
سؤال پژوهشی می‌تواند این باشد که آیا برآورد $\Delta(s)$ روی خودِ دستگاه، با بودجه شات محدود، می‌تواند زمان‌بندی را نسبت به جاروب خطی بهبود دهد، در حالی که هزینه اندازه‌گیری و خطای تروتر نیز در معیار وارد شوند؟

\paragraph{طیف‌سنجی مدل اسپینی روی کیوبیت‌های ابررسانا یا اتم خنثی.}
در این سناریو خروجی اصلی می‌تواند یک قله طیفی در تبدیل فوریه داده زمانی باشد.
سؤال این است که کدام مشاهده‌پذیر محلی گاف هدف را نشان می‌دهد و چگونه می‌توان اثر تقارن، انتخاب حالت اولیه و نویز را از گاف فیزیکی جدا کرد؟

\paragraph{شیمی و مواد در فضای فعال کوچک.}
گاف سینگلت--تریپلت یا شکافت دو حالت کم‌انرژی می‌تواند کمیتی شیمیایی یا فیزیکی باشد~\cite{McArdle2020,GapHiggott2019}.
در اینجا سؤال روش‌شناختی این است که آیا \lr{VQD}/\lr{QSE} روی یک فضای فعال با گواهی خطا، نسبت به \lr{CASCI} یا سایر روش‌های کلاسیک فقط بازتولید یک محاسبه کوچک است، یا در رژیمی که روش کلاسیک مناسب دشوار می‌شود مزیتی قابل‌اندازه‌گیری ایجاد می‌کند؟

\paragraph{حفاظت گاف در حافظه و گیت آدیاباتیک.}
گاف بزرگ‌تر می‌تواند نرخ برخی تحریک‌ها را کاهش دهد، ولی اجرای طولانی‌تر می‌تواند زمان مواجهه با ناهمدوسی و خطا را افزایش دهد.
بنابراین سؤال کمّی این نیست که «گاف بزرگ‌تر بهتر است یا نه»، بلکه این است که برای یک مدل نویزدار مشخص، آیا افزایش گاف و در نتیجه امکان اجرای آهسته‌تر، از نظر احتمال موفقیت نهایی بهتر از اجرای سریع‌تر با یک روش کنترل جبران‌کننده است؟

\section{سختی مسئله و آنچه نباید ادعا کرد}

تعیین انرژی پایه هامیلتونی محلی با دقت معکوس چندجمله‌ای در بدترین حالت \lr{QMA}-کامل است~\cite{GapGharibian2015}.
تصمیم‌گیری درباره ویژگی‌های گاف در حد ترمودینامیکی نیز برای برخی خانواده‌های ساختگی می‌تواند با پیچیدگی بسیار بالایی همراه باشد~\cite{GapCubitt2015}.

این نتایج مهم‌اند، ولی نباید از آنها نتیجه گرفت که هر مدل فیزیکی کوچک یا هر زنجیره آیزینگ مسئله‌ای سخت از دید محاسبات کلاسیک است.
برای اندازه‌های متناهی، قطری‌سازی مستقیم یا روش‌هایی مانند لانچوس و \lr{DMRG} ممکن است بسیار کارآمد باشند.
بنابراین در مقاله باید میان \emph{سختی بدترین حالت}، \emph{سختی خانواده مدل مورد مطالعه} و \emph{هزینه واقعی محاسبه برای اندازه‌های مورد نظر} تفاوت گذاشت.

همچنین $n$ کیوبیت به‌تنهایی معیار دشواری نیست.
آنچه در مطالعه مقیاس‌پذیری اهمیت دارد رفتار کمیت مورد نظر با افزایش اندازه سیستم است؛ برای گاف، این پرسش به شکل

\begin{equation}
	\Delta_{\min}(n)
\end{equation}

ظاهر می‌شود و رژیم چندجمله‌ای در برابر نمایی در بخش~\ref{sec:complexity-scaling} بررسی شده است.

از همین رو، ادعای «مزیت کوانتومی» نیز نباید تنها از مشاهده یک گاف کوچک یا اجرای یک مدار روی تعداد بیشتری کیوبیت نتیجه گرفته شود.
مقایسه باید هم‌زمان دقت، زمان، عمق مدار، تعداد شات، هزینه کنترل و بهترین مرجع کلاسیک در دسترس را در نظر بگیرد.

\section{مرور هدفمند منابع برای شروع پژوهش}

این فهرست جایگزین بخش~\ref{sec:gap-papers} نیست؛ آن بخش جزئیات مقالات ۲۰۲۵--۲۰۲۶ را برای مسیر آدیاباتیک دارد.
اینجا لنگرهای مروری و روش‌ساز آمده‌اند.

\begin{itemize}
	\item \textbf{شبیه‌سازی کوانتومی:}
	\lr{Georgescu--Ashhab--Nori} و چشم‌انداز شبیه‌سازی و مزیت عملی~\cite{Georgescu2014,Daley2022,CiracZoller2012}.
	
	```
	\item \textbf{آدیاباتیک و نقش گاف:}
	\lr{Farhi et al.}، مرور \lr{Albash--Lidar}، کران \lr{Jansen--Ruskai--Seiler} و زمان‌بندی \lr{Roland--Cerf}~\cite{Farhi2000,AlbashLidar2018,GapJansen2007,GapRoland2002}.
	
	\item \textbf{طیف کلاسیک:}
	مرور \lr{DMRG/MPS} و روش‌های طیفی کلاسیک~\cite{GapSchollwock2011,Cirac2021}.
	
	\item \textbf{روش‌های وردشی:}
	\lr{VQE}، مرور \lr{Tilly} و \lr{VQD}~\cite{Peruzzo2014,Tilly2022,GapHiggott2019}.
	
	\item \textbf{زیرفضا و کریلوف:}
	\lr{QSE}، تحلیل کریلوف، اجرای پردازنده و نمونه‌برداری کریلوف~\cite{McClean2017QSE,Kirby2023Krylov,KrylovNatComm2025,GapYu2025}.
	
	\item \textbf{مهندسی دینامیک و میان‌بر آدیاباتیک:}
	\lr{Berry}، \lr{Sels--Polkovnikov}، روش‌های مقابله با گاف کوچک و پیاده‌سازی دیجیتال--آنالوگ~\cite{Berry2009STA,GapSels2017,Gueneau2025CD,Kumar2024DACQO}.
	
	\item \textbf{برآورد گاف از سیگنال و در افق \lr{FTQC}:}
	\cite{GapLee2025,GapCugini2025,GapCarrera2026}.
	
	\item \textbf{سختی محاسباتی:}
	\cite{GapGharibian2015,GapCubitt2015}.
	```
	
\end{itemize}

\section{مسئله‌های باز و طرح مقاله}
\label{sec:gap-paper-tracks}

یک مقاله در این حوزه وقتی شانس دارد که \emph{یک ادعا}، \emph{یک خانواده مدل}، و \emph{یک معیار هزینه} را هم‌زمان محدود کند.
سؤال مادر را می‌توان به‌صورت زیر خلاصه کرد:

\begin{equation}
	\boxed{
		\text{چگونه می‌توان گاف را برآورد، کنترل و از آن بهره‌برداری کرد،
			با درنظرگرفتن محدودیت منابع کوانتومی؟}
	}
\end{equation}

از این دیدگاه، سه مسیر پژوهشی زیر مکمل یکدیگرند و با فعالیت گروه در شبیه‌سازی، آیزینگ، آدیاباتیک و \lr{Qiskit} سازگارند:

\begin{enumerate}
	\item \textbf{برآورد گاف مسیر با بودجه شات، سپس زمان‌بندی.}
	
	```
	فرضیه:
	روی زنجیره \lr{TFIM} با $n$ متوسط، تخمین نوسانی یا زیرفضا‌ییِ $\Delta(s)$ با خطای نسبی معلوم می‌تواند زمان‌بندی تطبیقی را به‌گونه‌ای تعیین کند که احتمال موفقیت نسبت به جاروب خطی، در زمان کل برابر، بهبود یابد.
	
	مرجع:
	قطری‌سازی یا فرمیون آزاد برای اندازه‌های کوچک~\cite{GapPfeuty1970}؛ پروتکل‌های اندازه‌گیری و پردازش سیگنال~\cite{GapLee2025,GapCugini2025}.
	
	خروجی مقاله:
	منحنی $P_{\mathrm{succ}}(T)$ در برابر خطای برآورد گاف و هزینه شات، نه فقط یک عدد $\Delta_{\min}$.
	
	\item \textbf{تقریب موضعی رانش جبران‌کننده در برابر تغییر مسیر.}
	
	فرضیه:
	برای همان خانواده، \lr{CD} محلی با شعاع محدود در رژیم گاف چندجمله‌ای می‌تواند سودمند باشد، ولی در رژیم گاف نمایی ممکن است به اشباع برسد؛ در برخی رژیم‌ها تغییر مسیر با هامیلتونی کاتالیزگر می‌تواند گزینه بهتری باشد~\cite{GapSels2017,Gueneau2025CD}.
	
	معیار:
	وفاداری نهایی در برابر هزینه کنترل، عمق مدار یا تعداد گیت برابر.
	در این مسیر هدف فقط افزایش $\Delta_{\min}$ نیست، بلکه مقایسه هزینه واقعی دو راه برای کاهش خطای نهایی است.
	
	\item \textbf{گواهی گاف در زیرفضای کریلوف، نه فقط تفاضل وردشی.}
	
	فرضیه:
	با کنترل باقیمانده و عدد وضعیت $S$، می‌توان بازه‌ای مانند
	$[\Delta_{\min}^{\mathrm{lo}},\Delta_{\min}^{\mathrm{hi}}]$
	برای یک مدل اسپینی دوبعدی کوچک به‌صورت قابل‌گزارش به‌دست آورد، در حالی که تفاضل دو تخمین \lr{VQE} به‌تنهایی چنین بازه‌ای را تضمین نمی‌کند~\cite{Kirby2023Krylov,KrylovNatComm2025,McClean2017QSE}.
	
	این مسیر به \lr{early-FTQC} نزدیک‌تر است و از صرفاً بازتولید \lr{VQD} فاصله می‌گیرد.
	```
	
\end{enumerate}

مسیرهای~\ref{sec:gap-projects} را می‌توان زیرمجموعه‌ای از همین سه ادعا دید.
دو اسپین برای راستی‌آزمایی، بررسی مفهومی و آزمون اولیه بسیار مفید است، اما به‌تنهایی شواهدی برای مقیاس‌پذیری یا مقاله روش‌شناختی فراهم نمی‌کند.

\subsection{چک‌لیست قبل از شروع محاسبه}

پیش از اجرای محاسبات اصلی، باید دست‌کم پرسش‌های زیر پاسخ داده شوند:

\begin{enumerate}
	\item کدام گاف مورد نظر است: گاف طیفی کل، گاف یک بخش تقارنی، گاف قابل‌مشاهده، یا گاف مسیر $\Delta_{\min}(s)$؟
	\item واحد انرژی چیست و آیا مقیاس انرژی با افزایش $n$ ثابت نگه داشته می‌شود؟
	\item اگر مسیر داریم، آیا $\Delta(s)$ را در کل مسیر بررسی کرده‌ایم یا فقط چند نقطه را؟
	\item مرجع کلاسیک تا چه $n$ می‌تواند جواب دقیق یا کنترل‌شده بدهد؟
	\item هزینه را با چه واحدی می‌سنجیم: زمان فیزیکی، عمق تروتر، تعداد شات، تعداد گیت، یا پالس آنالوگ؟
	\item آیا نویز و ناهمدوسی مدل شده‌اند یا فقط سامانه بسته مطالعه می‌شود؟
	\item ادعای اصلی بهبود $\Delta$ است یا بهبود $P_{\mathrm{succ}}$ در منابع برابر؟
	\item اگر گاف افزایش یافته است، آیا هزینه هامیلتونی یا کنترل جدید نیز در مقایسه وارد شده است؟
	\item اگر گاف از داده تجربی استخراج شده است، عدم‌قطعیت تخمین آن چگونه گزارش می‌شود؟
\end{enumerate}

\section{ارتباط}

تعریف گاف و نقش تقارن در بخش~\ref{sec:gap-research}؛ از گاف تا زمان در~\ref{sec:gap-runtime}؛ مقیاس‌بندی در~\ref{sec:complexity-scaling}؛ روش‌های عددی مسیر دو اسپین در~\ref{sec:gap-methods}--\ref{sec:gap-projects}.
ابزار دیجیتال در فصل~\ref{chap:digital-tools}؛ \lr{QPE}/\lr{VQE}/\lr{QAOA} در فصل~\ref{chap:spectral-tools}؛ بستر آنالوگ در فصل~\ref{chap:hardware}.
مدل آیزینگ در بخش~\ref{sec:ising}.

ارتباط اصلی با فصل~\ref{chap:evo-adi} نیز از این جهت است که در آنجا گاف در چارچوب یک مسیر آدیاباتیک و از منظر زمان اجرا و احتمال موفقیت بررسی شد، در حالی که در این فصل گاف به‌عنوان یک موضوع مستقل برای برآورد، تفسیر و مهندسی دنبال می‌شود.
به این ترتیب، مسیر آدیاباتیک یک حالت خاص از کاربرد گاف در فصل حاضر است، نه تعریف‌کننده کل دامنه آن.

\section{چه وقت استفاده می‌شود}

این فصل زمانی نقطه شروع مناسبی است که سؤال پژوهش یکی از این موارد باشد:

\begin{itemize}
	\item گاف یک مدل چیست و چگونه می‌توان آن را با دقت مشخص برآورد کرد؟
	\item چگونه می‌توان گاف را از داده دینامیکی یا اندازه‌گیری روی سخت‌افزار استخراج کرد؟
	\item چرا گاف کوچک در یک مسیر یا مدل ظاهر می‌شود؟
	\item آیا می‌توان با تغییر مسیر، رمزگذاری یا کنترل، اثر گاف کوچک را کاهش داد؟
	\item آیا یک روش کوانتومی برای برآورد یا مهندسی گاف در برابر بهترین روش کلاسیک مزیتی دارد؟
\end{itemize}

اگر فقط می‌خواهید یک مسیر آدیاباتیک مشخص را بفهمید و اثر گاف آن بر زمان اجرا را بررسی کنید، فصل~\ref{chap:evo-adi} کافی است؛ این فصل برای تبدیل «گاف» از یک کمیت کمکی به یک خط پژوهشی مستقل طراحی شده است.

\begin{remark}[جای توسعه]
	یک بنچمارک مشترک روی \lr{TFIM} باز و یک مدل با گاف نمایی، با همان بودجه شات؛ اتصال گاف لیوویلی به نرخ موفقیت در معادله مادر؛ مقایسه صریح \lr{QAOA} با مسیر آدیاباتیک از نظر گاف مؤثر، هزینه کنترل و احتمال موفقیت، نه فقط مقدار نهایی تابع هزینه.
\end{remark}
